\section{Introduction}
\label{sec:introduction}

A model asked to charge a credit card might emit
\texttt{stripe.Charge.create(...)} --- deprecated since 2019 --- or
\texttt{stripe.PaymentIntent.create(...)}, the current API. Both arrive in
the developer's editor with the same syntactic fluency, the same confident
indentation, and no visible marker distinguishing a well-chosen call from
one the model guessed at. At the method-selector token, the two choices
split the probability mass near 50/50 ($0.52$ vs.\ $0.47$); the output
exposes none of this. Failures of AI code generation cluster at exactly
these points: version-split APIs, renamed SDK methods, signature changes
between library releases. The model is uncertain; the output is not.

\paragraph{Intra- versus inter-generational uncertainty.}
Existing uncertainty-quantification methods for LLMs are overwhelmingly
\emph{inter-generational}: they sample the same prompt multiple times and
measure disagreement across completions~\cite{kuhn2023semantic,uqlm2024}.
This works when the model's uncertainty surfaces as instability across
re-rolls --- factual questions, open-ended reasoning. Code is not such a
task. A model with a $0.52$ prior on the deprecated \texttt{Charge} API
samples that branch nearly every time; once the early tokens commit, the
rest of the function is forced by context. The uncertainty exists, but it
lives \emph{inside a single generation}, in the token distribution at each
decoding step, not \emph{across} generations. We call this the
\emph{intra- versus inter-generational} distinction. It is the conceptual
foundation of this paper: UQLM, the closest published tool for LLM
uncertainty quantification, achieves $5\%$ detection on our benchmark
where \eps achieves $85\%$ (\S\ref{sec:comparison}).

\paragraph{\eps: a recall-maximizing filter.}
We introduce \eps (epsilon), a per-token normalized Shannon entropy over
the top-$k$ log-probabilities returned during generation. After AST-based
filtering to keep only semantically meaningful tokens (attribute accesses,
keyword arguments, import targets --- not whitespace, comments, or
cosmetic identifiers), we aggregate via max-pooling to obtain a single
score per function. \eps is not a detector --- it identifies where the
model was uncertain at decoding time. At a threshold of $\eps \ge 0.30$
it catches $100\%$ of true failures on our benchmark, at the cost of
roughly $73\%$ false alarms.

\paragraph{The two-stage system.}
A $27\%$-precision filter is not meant to be used in isolation. Stage 2
is a parallel LLM review loop: every flagged entry is dispatched, in
parallel, to a reviewer model that returns \textsc{cleared} or
\textsc{confirmed}. At current API prices this costs fractions of a cent
per entry and completes in seconds. \eps has perfect recall but poor
precision; the reviewer has good precision but would be too expensive per
line. Filter plus reviewer gives a system with near-$100\%$ end-to-end
precision while missing nothing the filter catches. AI-assisted developers
now write code several times faster than they can review it; a parallel
review system scales with generation speed and moves review off the
developer's critical path.

\paragraph{Contributions.}
\begin{enumerate}[noitemsep]
    \item The intra/inter-generational distinction, with empirical
    confirmation: $5\%$ vs.\ $85\%$ detection for UQLM vs.\ \eps on the
    same benchmark.
    \item \eps: an AST-filtered, max-aggregated token-level entropy score,
    with a three-type uncertainty taxonomy (cosmetic, consequential,
    implementation-approach) that motivates the filter design.
    \item A two-stage filter-plus-review architecture, with a context-%
    accumulation variant, evaluated on 201 entries across three models:
    near-$100\%$ end-to-end precision at $100\%$ recall.
\end{enumerate}

Implementation and benchmark are at
\url{https://github.com/michaeljaquith/epsilon-code}.
